\chapter{GPT 모델 완성과 학습 루프}
\label{ch:gpt_training}

이 장에서는 6장의 트랜스포머 블록을 $N$층 쌓아 올려 완전한 \keyterm{GPT 모델}을 만들고, 이 모델을 학습하는 루프를 구축한다. 데이터셋, 손실 함수, 옵티마이저, 스케줄러를 모두 직접 구현한다.

\section{GPT 모델 구조}

GPT는 트랜스포머 디코더를 $N$층 쌓은 자기회귀 언어 모델이다. 전체 구조는 다음과 같다:

\begin{enumerate}
  \item 토큰 ID $\to$ [Token + Position Embedding]
  \item [Transformer Block] $\times N$
  \item [Final LayerNorm]
  \item [LM Head: $d_{model} \to V$]
  \item logits
\end{enumerate}

\begin{figure}[htbp]
\centering
\begin{tikzpicture}[
  every node/.style={font=\small\sffamily},
  box/.style={draw=inkblue, fill=inkblue!8, rounded corners=2pt, minimum width=4cm, minimum height=0.7cm},
  rep/.style={draw=inkgray, fill=inkgray!10, rounded corners=2pt, minimum width=4cm, minimum height=0.7cm, dashed}
]
  \node[box] (in) at (0, 0) {token IDs [batch, seq]};
  \node[box, below=0.4cm of in] (emb) {Embedding (token + position)};
  \node[rep, below=0.4cm of emb] (blk) {Transformer Block $\times N$};
  \node[box, below=0.4cm of blk] (ln) {Final LayerNorm};
  \node[box, below=0.4cm of ln] (head) {LM Head ($d_{model} \to V$)};
  \node[box, below=0.4cm of head] (out) {logits [batch, seq, V]};

  \draw[-{Stealth[length=4pt]}, inkblue] (in) -- (emb);
  \draw[-{Stealth[length=4pt]}, inkblue] (emb) -- (blk);
  \draw[-{Stealth[length=4pt]}, inkblue] (blk) -- (ln);
  \draw[-{Stealth[length=4pt]}, inkblue] (ln) -- (head);
  \draw[-{Stealth[length=4pt]}, inkblue] (head) -- (out);

  % weight tying 표시
  \draw[inkred, dashed, <->] (emb.east) to[bend left=30] node[right, inkred, font=\scriptsize] {weight tying} (head.east);
\end{tikzpicture}
\caption{GPT 모델 구조. 임베딩 가중치와 LM Head 가중치를 공유(weight tying)하여 파라미터 절약.}
\label{fig:gpt}
\end{figure}

\section{가중치 타이 (Weight Tying)}

임베딩 행렬 $E$ ($V \times d_{model}$)과 LM Head 행렬 $W_{out}$ ($d_{model} \times V$)은 사실 같은 정보를 담고 있다. 토큰 $i$의 임베딩은 $E[i]$이고, 로짓 계산 시 $W_{out}^T[i]$가 사용된다. 이 둘을 공유하면 파라미터 수를 $V \times d_{model}$만큼 절약할 수 있다.

\begin{lstlisting}[language=Python]
class GPT(nn.Module):
    def __init__(self, config: ModelConfig):
        super().__init__()
        self.config = config
        self.embedding = EmbeddingStack(
            vocab_size=config.vocab_size,
            d_model=config.d_model,
            max_len=config.context_length,
            pad_id=config.pad_token_id,
        )
        self.blocks = nn.ModuleList([
            TransformerBlock(config, layer_norm_style="pre")
            for _ in range(config.n_layers)
        ])
        self.ln_f = nn.LayerNorm(config.d_model, eps=config.layer_norm_eps)
        self.lm_head = nn.Linear(config.d_model, config.vocab_size, bias=False)
        # 가중치 타이
        self.lm_head.weight = self.embedding.token_emb.embedding.weight
\end{lstlisting}

\begin{notebox}[가중치 타이의 장단점]
\textbf{장점}: 파라미터 절약 (GPT-2 small 기준 약 1/3 감소), 임베딩과 출력의 일관성.

\textbf{단점}: 임베딩이 출력 레이어의 그래디언트도 받아 학습이 더 복잡해짐. 일부 모델(LLaMA 등)은 타이를 사용하지 않고 별도 학습.

GPT-2, GPT-3는 타이 사용. LLaMA, Mistral은 미사용. 이 책에서는 GPT 스타일을 따라 타이를 사용한다.
\end{notebox}

\section{Forward Pass}

\begin{lstlisting}[language=Python]
    def forward(self, token_ids, return_hidden=False):
        batch, seq = token_ids.shape
        # 인과 마스크 생성
        mask = make_causal_mask(seq, device=token_ids.device)
        # 임베딩
        x = self.embedding(token_ids)
        # 트랜스포머 블록 통과
        for block in self.blocks:
            x = block(x, mask=mask)
        # 최종 정규화 (Pre-LN이므로 마지막에 필요)
        x = self.ln_f(x)
        # 로짓
        logits = self.lm_head(x)
        if return_hidden:
            return logits, x
        return logits
\end{lstlisting}

\begin{infobox}[왜 마지막에 LayerNorm이 필요한가?]
Pre-LN 블록은 각 블록의 입력을 정규화한다. 하지만 마지막 블록의 출력은 정규화되지 않은 상태로 잔차 스트림에 더해진다. 이것을 LM Head에 넣기 전에 정규화해야 출력 분포가 안정적이다. 그래서 \code{ln\_f} (final LayerNorm)가 필요하다.

Post-LN에서는 각 블록 출력이 이미 정규화되어 있으므로 \code{ln\_f}가 불필요하다.
\end{infobox}

\section{데이터셋: 슬라이딩 윈도우}

언어 모델링 데이터셋은 텍스트를 토큰화한 뒤, 고정 길이 청크로 분할한다. 각 청크에서 input은 처음 $L-1$ 토큰, target은 마지막 $L-1$ 토큰 (한 칸 shift).

\begin{lstlisting}[language=Python]
class LMDataset(Dataset):
    def __init__(self, text, tokenizer, context_length=128,
                 add_bos=True, add_eos=True):
        self.context_length = context_length
        # 텍스트를 토큰화하여 긴 시퀀스로
        if isinstance(text, str):
            text = [text]
        all_ids = []
        for t in text:
            ids = tokenizer.encode(t, add_bos=add_bos, add_eos=add_eos)
            all_ids.extend(ids)
        self._ids = all_ids
        # context_length+1 크기 청크로 분할
        cl = context_length + 1
        self._chunks = []
        for i in range(0, len(all_ids), cl):
            chunk = all_ids[i:i+cl]
            if len(chunk) < cl:
                chunk = chunk + [tokenizer.special.pad_id] * (cl - len(chunk))
            self._chunks.append(chunk)

    def __getitem__(self, idx):
        chunk = self._chunks[idx]
        input_ids = torch.tensor(chunk[:-1], dtype=torch.long)
        target_ids = torch.tensor(chunk[1:], dtype=torch.long)
        return input_ids, target_ids
\end{lstlisting}

\begin{figure}[htbp]
\centering
\begin{tikzpicture}[
  every node/.style={font=\small\sffamily},
  tok/.style={draw=inkblue, fill=inkblue!15, minimum width=0.5cm, minimum height=0.5cm},
  dstin/.style={draw=inkblue, fill=inkblue!30, minimum width=0.5cm, minimum height=0.5cm},
  tgt/.style={draw=inkred, fill=inkred!30, minimum width=0.5cm, minimum height=0.5cm}
]
  % 시퀀스
  \node[inkgray, anchor=east] at (-0.3, 0) {tokens:};
  \foreach \i/\t in {0/B, 1/t, 2/h, 3/e, 4/c, 5/a, 6/t, 7/E} {
    \node[tok] at (\i*0.55, 0) {\t};
  }

  \node[inkgray, anchor=east] at (-0.3, -0.8) {input:};
  \foreach \i/\t in {0/B, 1/t, 2/h, 3/e, 4/c, 5/a, 6/t} {
    \node[dstin] at (\i*0.55, -0.8) {\t};
  }

  \node[inkgray, anchor=east] at (-0.3, -1.6) {target:};
  \foreach \i/\t in {0/t, 1/h, 2/e, 3/c, 4/a, 5/t, 6/E} {
    \node[tgt] at (\i*0.55, -1.6) {\t};
  }

  % 화살표
  \foreach \i in {0, 1, 2, 3, 4, 5, 6} {
    \draw[-{Stealth[length=3pt]}, inkgray!50] (\i*0.55, -0.3) -- (\i*0.55, -0.55);
  }

  \node[inkgray, font=\scriptsize, anchor=west] at (5, -0.8) {input은 target보다 한 칸 앞};
  \node[inkgray, font=\scriptsize, anchor=west] at (5, -1.6) {B=<bos>, E=<eos>};
\end{tikzpicture}
\caption{슬라이딩 윈도우 데이터셋. input과 target은 같은 시퀀스를 한 칸 shift한 것. 모델은 input을 보고 target을 예측.}
\label{fig:sliding-window}
\end{figure}

\section{손실 함수: Cross-Entropy}

언어 모델의 손실은 \keyterm{교차 엔트로피(cross-entropy)}다. 모델이 예측한 확률 분포와 정답(원-핫) 사이의 차이를 측정한다.

\begin{formula}[Cross-Entropy Loss]
\[
\mathcal{L} = -\frac{1}{N} \sum_{i=1}^{N} \log P(y_i \mid x_{<i})
\]
$y_i$는 정답 토큰, $P$는 모델의 예측 확률. 낮을수록 좋다.
\end{formula}

\begin{lstlisting}[language=Python]
def cross_entropy_loss(logits, targets, ignore_index=-100):
    # logits: [batch, seq, V], targets: [batch, seq]
    batch, seq, vocab = logits.shape
    logits_flat = logits.view(-1, vocab)
    targets_flat = targets.view(-1)
    return F.cross_entropy(logits_flat, targets_flat, ignore_index=ignore_index)
\end{lstlisting}

\subsection{Label Smoothing (선택)}

정답에 100\% 확신하는 대신 약간의 불확실성을 허용하여 일반화 성능을 높이는 기법.

\begin{lstlisting}[language=Python]
def label_smoothing_loss(logits, targets, smoothing=0.1, ignore_index=-100):
    """정답 토큰에 1-smoothing 확률, 나머지에 smoothing/vocab 확률."""
    log_probs = F.log_softmax(logits.view(-1, logits.size(-1)), dim=-1)
    targets_flat = targets.view(-1)
    nll_loss = -log_probs.gather(1, targets_flat.clamp(min=0).unsqueeze(1)).squeeze(1)
    smooth_loss = -log_probs.mean(dim=-1)
    mask = (targets_flat != ignore_index).float()
    nll_loss = (nll_loss * mask).sum() / mask.sum().clamp(min=1)
    smooth_loss = (smooth_loss * mask).sum() / mask.sum().clamp(min=1)
    return (1.0 - smoothing) * nll_loss + smoothing * smooth_loss
\end{lstlisting}

\section{옵티마이저: AdamW}

\keyterm{AdamW}는 Adam에 weight decay를 결합한 최적화 알고리즘이다. 핵심은 바이어스와 LayerNorm에는 weight decay를 적용하지 않는 것.

\begin{lstlisting}[language=Python]
def build_optimizer(model, config):
    # 파라미터 그룹 분리: decay / no_decay
    decay, no_decay = [], []
    for name, p in model.named_parameters():
        if not p.requires_grad: continue
        if p.ndim < 2 or "bias" in name or "ln" in name.lower():
            no_decay.append(p)
        else:
            decay.append(p)
    return AdamW(
        [{"params": decay, "weight_decay": config.weight_decay},
         {"params": no_decay, "weight_decay": 0.0}],
        lr=config.learning_rate, betas=(0.9, 0.95)
    )
\end{lstlisting}

\begin{notebox}[왜 바이어스와 LayerNorm은 weight decay에서 빼는가?]
Weight decay는 가중치를 0으로 끌어당기는 정규화다. 행렬 가중치에는 과적합 방지 효과가 있지만, 바이어스와 LayerNorm의 $\gamma, \beta$는 차원이 낮아 weight decay의 효과가 미미하고 오히려 학습을 방해할 수 있다. PyTorch 커뮤니티의 표준 관행이다.
\end{notebox}

\section{학습률 스케줄러}

학습 초기에는 학습률을 서서히 올리고(warmup), 이후 코사인 곡선으로 감소시킨다. 이 \keyterm{warmup + cosine} 스케줄은 GPT-3에서 표준으로 사용된다.

\begin{lstlisting}[language=Python]
def build_scheduler(optimizer, config, total_steps):
    warmup = config.warmup_steps
    def lr_lambda(step):
        if step < warmup:
            return float(step) / max(1, warmup)
        progress = (step - warmup) / max(1, total_steps - warmup)
        if config.lr_scheduler == "cosine":
            return 0.5 * (1.0 + math.cos(math.pi * min(progress, 1.0)))
        return 1.0
    return LambdaLR(optimizer, lr_lambda)
\end{lstlisting}

\begin{figure}[htbp]
\centering
\begin{tikzpicture}[
  scale=0.9,
  every node/.style={font=\small\sffamily}
]
  % 축
  \draw[inkblue, line width=0.5pt, ->] (0, 0) -- (11, 0) node[right] {step};
  \draw[inkblue, line width=0.5pt, ->] (0, 0) -- (0, 3.5) node[above] {LR};
  % warmup 구간 (선형 증가)
  \draw[inkred, line width=1pt] (0, 0) -- (2, 2.5);
  \node[inkred, font=\scriptsize, anchor=south] at (1, 2.6) {warmup};
  % cosine 구간
  \draw[inkblue, line width=1pt, domain=2:10, samples=50] plot (\x, {2.5*0.5*(1+cos(180*(\x-2)/8))});
  \node[inkblue, font=\scriptsize, anchor=south] at (6, 2.6) {cosine decay};
  % 최대 LR 표시
  \draw[inkgray, dashed] (0, 2.5) -- (2, 2.5);
  \node[inkgray, font=\scriptsize, anchor=east] at (-0.1, 2.5) {$lr_{max}$};
\end{tikzpicture}
\caption{Warmup + Cosine 학습률 스케줄. 처음에는 선형으로 올리고, 이후 코사인 곡선으로 감소.}
\label{fig:lr-schedule}
\end{figure}

\subsection{왜 warmup이 필요한가?}

학습 초기에는 가중치가 무작위에 가깝다. 큰 학습률을 바로 적용하면 그래디언트가 불안정하여 발산할 수 있다. 작은 학습률로 시작하여 점진적으로 올리면 모델이 안정적인 영역으로 들어간다.

특히 어텐션의 초기 그래디언트는 매우 클 수 있는데 (softmax 포화 등), warmup이 이를 완화한다. Pre-LN 구조는 Post-LN보다 warmup에 덜 민감하지만, 여전히 사용하는 것이 안전하다.

\section{그래디언트 클리핑}

그래디언트가 너무 커서 가중치가 크게 변하는 것을 방지하기 위해 그래디언트 norm을 제한한다.

\begin{lstlisting}[language=Python]
if config.grad_clip > 0:
    torch.nn.utils.clip_grad_norm_(model.parameters(), config.grad_clip)
\end{lstlisting}

보통 \code{grad\_clip = 1.0}을 사용한다. 그래디언트 norm이 1.0을 초과하면 1.0으로 스케일링한다.

\section{트레이너}

모든 컴포넌트를 묶어 학습 루프를 만든다.

\begin{lstlisting}[language=Python]
class Trainer:
    def __init__(self, model, train_dataset, config, model_config, eval_dataset=None):
        self.model = model
        self.config = config
        self.device = torch.device(config.device)
        self.model.to(self.device)
        self.train_loader = DataLoader(train_dataset, batch_size=config.batch_size,
                                       shuffle=True, collate_fn=collate_batch)
        # 총 스텝 수 계산
        total_steps = len(train_dataset) // config.batch_size * config.max_epochs
        self.optimizer = build_optimizer(model, config)
        self.scheduler = build_scheduler(self.optimizer, config, total_steps)
        # ...

    def train(self):
        self.model.train()
        for inputs, targets in self.train_loader:
            inputs, targets = inputs.to(self.device), targets.to(self.device)
            logits = self.model(inputs)
            loss = cross_entropy_loss(logits, targets,
                                       ignore_index=self.model_config.pad_token_id)
            self.optimizer.zero_grad(set_to_none=True)
            loss.backward()
            if self.config.grad_clip > 0:
                torch.nn.utils.clip_grad_norm_(self.model.parameters(),
                                                self.config.grad_clip)
            self.optimizer.step()
            self.scheduler.step()
\end{lstlisting}

\section{Perplexity}

\keyterm{퍼플렉서티(Perplexity, PPL)}는 언어 모델의 대표적 평가 지표다.

\begin{formula}[Perplexity]
\[
\text{PPL} = \exp\left(\frac{1}{N}\sum_i \log P(y_i | x_{<i})\right) = \exp(\text{loss})
\]
PPL은 모델이 다음 토큰을 예측할 때 ``평균적으로 몇 개의 후보로 헤매는가''를 나타낸다. PPL=1이면 완벽, PPL=10이면 10개 후보로 좁힐 수 있다는 뜻.
\end{formula}

\begin{infobox}[PPL 해석]
\begin{itemize}
  \item PPL = 1: 모델이 다음 토큰을 100\% 확신. 이론적 최소.
  \item PPL = 10: 10개 토큰 중 하나로 좁힘. 언어 모델로서 쓸만한 수준.
  \item PPL = 100: 100개 토큰 중 하나. 거의 무작위.
  \item PPL = $V$ (어휘 크기): 완전 무작위. 학습 안 됨.
\end{itemize}
GPT-2 small의 WikiText-103 PPL은 약 17. GPT-3 175B는 약 11.
\end{infobox}

\section{체크포인트 저장/로드}

학습 중 모델 가중치를 주기적으로 저장하여, 중단 후 재개하거나 추론에 사용할 수 있다.

\begin{lstlisting}[language=Python]
def save_checkpoint(self, filename):
    path = self.out_dir / filename
    torch.save({
        "model_state": self.model.state_dict(),
        "optimizer_state": self.optimizer.state_dict(),
        "scheduler_state": self.scheduler.state_dict(),
        "global_step": self.global_step,
        "model_config": self.model_config.to_dict(),
        "trainer_config": self.config.to_dict(),
    }, path)

def load_checkpoint(self, path):
    ckpt = torch.load(path, map_location=self.device)
    self.model.load_state_dict(ckpt["model_state"])
    self.optimizer.load_state_dict(ckpt["optimizer_state"])
    self.scheduler.load_state_dict(ckpt["scheduler_state"])
    self.global_step = ckpt["global_step"]
\end{lstlisting}

\section{요약}

\begin{itemize}
  \item GPT = Embedding + Transformer Block $\times N$ + LayerNorm + LM Head.
  \item 가중치 타이로 임베딩과 LM Head 가중치를 공유하여 파라미터 절약.
  \item Pre-LN 구조에서는 마지막에 \code{ln\_f}가 필요.
  \item 언어 모델링은 슬라이딩 윈도우로 (input, target) 쌍 생성.
  \item 손실은 cross-entropy, 옵티마이저는 AdamW (bias/LayerNorm은 weight decay 제외).
  \item Warmup + Cosine 학습률 스케줄 + 그래디언트 클리핑이 표준.
  \item Perplexity = $\exp(\text{loss})$로 모델 품질 평가.
\end{itemize}

\section{연습 문제}

\begin{enumerate}
  \item 가중치 타이를 사용하지 않으면 GPT-2 small (V=50257, d=768)에서 파라미터가 얼마나 늘어나는가?
  \item Warmup 스텝 수를 0으로 하면 학습이 어떻게 달라지는가?
  \item Label smoothing을 0.3으로 강하게 적용하면 어떤 효과가 있는가?
  \item 그래디언트 클리핑 없이 학습하면 어떤 현상이 발생할 수 있는가?
  \item PPL이 1.0인 모델이 실제로 존재할 수 있는가? 왜 그런가?
\end{enumerate}

다음 장에서는 학습된 모델로 텍스트를 생성하는 방법을 다룬다. Greedy, Top-K, Top-P 샘플링 전부 구현한다.
